% !TeX program = xelatex
% ============================================================================
%  第 13 课　它有什么用
% ----------------------------------------------------------------------------
%  编译        XeLaTeX，编译两遍
%  样式        ../../00-共享/preamble.tex
%  说明        本课不发单页
%  图片        一律 \imgph 占位，不要在这里手画 tikz
% ============================================================================

\documentclass[10pt]{beamer}
% ============================================================================
%  讲义样式文件 —— 字体 / 配色 / 校徽栏 / 顶部导航栏 / 页脚 / 常用方框
%  所有个性化设置都集中在这里，改一次即可用于所有讲义
%
%  用法：主文件里写
%      \documentclass[10pt]{beamer}
%      \input{preamble.tex}
%  依赖：同目录下的 logo_white.png（顶部校徽）
%  编译：必须用 XeLaTeX，而且必须编译两遍
%        （顶部导航栏的节名与小圆点记录在 .nav 文件里，只编一遍不会更新）
% ============================================================================
\usepackage{ctex}  % 支持中文
\usepackage{amsmath, amssymb}  % 数学公式支持
\usepackage{fontspec}  % 字体设置
\usepackage{unicode-math}  % 支持Unicode数学字体
\usepackage{xcolor}  % 颜色
\usepackage{tikz}  % 用于绘制页脚色条
\usepackage{booktabs}  % 表格

% 图片搜索路径 ------------------------------------------------------------------
% 校本课的资料分目录存放，本文件可能被 02-课件/第XX课/ 、03-讲义/ 下的主文件调用，
% 这里把几种可能的相对位置都列上，logo_white.png 放在 00-共享/ 即可被找到。
% 最后一项 00-共享/图片/ 是给 \imgfig 用的插图目录（图与课件不在同一层，靠它找到）。
\graphicspath{{./}{./图片/}{../../00-共享/}{../00-共享/}{00-共享/}{../../00-共享/图片/}{../00-共享/图片/}{00-共享/图片/}}

% 设置字体
\setCJKmainfont{SimSun}  % 设置中文字体为宋体
\setmainfont{Palatino Linotype}  % 设置英文字体为Palatino Linotype
\setmathfont{Latin Modern Math}  % 设置数学字体为Latin Modern Math
% 注意：beamer 默认把正文字体族设成无衬线（sans），Windows 下对应的中文字体是
% 微软雅黑，所以上面这行 SimSun 目前对正文不起作用。想让正文变成宋体，取消
% 下面这行的注释即可（标题栏、页脚、方框标题不受影响）。
% \renewcommand{\familydefault}{\rmdefault}

% 定义颜色 - 均为绿色系
\definecolor{maingreen}{RGB}{34,139,34}  % 森林绿（导航栏底色）
\definecolor{lightgreen}{RGB}{240,255,240}  % 蜜瓜白（帧标题栏底色）
\definecolor{accentgreen}{RGB}{46,139,87}  % 海绿（强调色）
\definecolor{textgreen}{RGB}{0,128,0}  % 纯绿（标题文字）
\definecolor{darkgreen}{RGB}{20,90,20}  % 深绿（顶部校徽栏底色）

% 设置Beamer主题
\usetheme{default}
\usecolortheme{default}

% 自定义主题
\setbeamercolor{background canvas}{bg=white}
\setbeamercolor{frametitle}{fg=textgreen, bg=lightgreen}
\setbeamercolor{title}{fg=textgreen}
\setbeamercolor{subtitle}{fg=textgreen}
\setbeamercolor{section in toc}{fg=accentgreen}
\setbeamercolor{normal text}{fg=black}
\setbeamercolor{structure}{fg=accentgreen}
\setbeamercolor{itemize item}{fg=accentgreen}
\setbeamercolor{itemize subitem}{fg=accentgreen}
\setbeamercolor{itemize subsubitem}{fg=accentgreen}
\setbeamercolor{block title}{fg=white, bg=accentgreen}
\setbeamercolor{block body}{bg=lightgreen}
% example 环境（底层是 exampleblock）用的是独立的一组颜色名，default 配色主题下
% 它们的底色是空的，所以只设 block 的话“例”会变成一行没有框的标题，这里接过来
\setbeamercolor{block title example}{parent=block title}
\setbeamercolor{block body example}{parent=block body}

% ==================== 常用定理类方框（五类配色）====================
% 五类各自的标题条底色，正文底色由同色自动调浅（!8!white）
\definecolor{boxThm}{RGB}{20,90,20}    % 一、定理类：深绿
\definecolor{boxDef}{RGB}{46,139,87}   % 二、定义类：海绿
\definecolor{boxExm}{RGB}{34,139,34}   % 三、例题类：森林绿
\definecolor{boxPrf}{RGB}{85,107,47}   % 四、过程类：橄榄绿
\definecolor{boxNote}{RGB}{0,121,107}  % 五、注记类：青绿

% 方框本体：先套本类配色，再调用 beamer 的 block（改动只在环境内部生效）
\newcommand{\gbox}[2]{%
  \setbeamercolor{block title}{fg=white,bg=#1}%
  \setbeamercolor{block body}{bg=#1!8!white}%
  \begin{block}{#2}}
\newcommand{\gboxend}{\end{block}}
% 标题后的可选补充名：\begin{theorem}[勾股定理] → 定理（勾股定理）
\newcommand{\gsub}[1]{\ifx\relax#1\relax\else（#1）\fi}


% 一、定理类（定理/推论/性质/命题/引理）
\renewenvironment{theorem}[1][]{\gbox{boxThm}{定理\gsub{#1}}}{\gboxend}
\renewenvironment{corollary}[1][]{\gbox{boxThm}{推论\gsub{#1}}}{\gboxend}
\newenvironment{character}[1][]{\gbox{boxThm}{性质\gsub{#1}}}{\gboxend}
\newenvironment{proposition}[1][]{\gbox{boxThm}{命题\gsub{#1}}}{\gboxend}
\renewenvironment{lemma}[1][]{\gbox{boxThm}{引理\gsub{#1}}}{\gboxend}

% 二、定义类（定义/公理）
\renewenvironment{definition}[1][]{\gbox{boxDef}{定义\gsub{#1}}}{\gboxend}
\newenvironment{axiom}[1][]{\gbox{boxDef}{公理\gsub{#1}}}{\gboxend}

% 三、例题类（例/练习/习题）
\renewenvironment{example}[1][]{\gbox{boxExm}{例\gsub{#1}}}{\gboxend}
\newenvironment{exercise}[1][]{\gbox{boxExm}{练习\gsub{#1}}}{\gboxend}
\newenvironment{homework}[1][]{\gbox{boxExm}{习题\gsub{#1}}}{\gboxend}

% 四、过程类（证明/解/法1~法4）
\renewenvironment{proof}[1][]{\gbox{boxPrf}{证明\gsub{#1}}}{\qedsymbol\gboxend}
\newenvironment{jie}[1][]{\gbox{boxPrf}{解\gsub{#1}}}{\gboxend}
\newenvironment{way1}[1][]{\gbox{boxPrf}{法1\gsub{#1}}}{\gboxend}
\newenvironment{way2}[1][]{\gbox{boxPrf}{法2\gsub{#1}}}{\gboxend}
\newenvironment{way3}[1][]{\gbox{boxPrf}{法3\gsub{#1}}}{\gboxend}
\newenvironment{way4}[1][]{\gbox{boxPrf}{法4\gsub{#1}}}{\gboxend}

% 五、注记类（经验总结/注/值得注意的是/分析/提示）
\newenvironment{experience}[1][]{\gbox{boxNote}{经验总结\gsub{#1}}}{\gboxend}
\newenvironment{zhu}[1][]{\gbox{boxNote}{注\gsub{#1}}}{\gboxend}
\newenvironment{att}[1][]{\gbox{boxNote}{值得注意的是\gsub{#1}}}{\gboxend}
\newenvironment{analyse}[1][]{\gbox{boxNote}{分析\gsub{#1}}}{\gboxend}
\newenvironment{hint}[1][]{\gbox{boxNote}{提示\gsub{#1}}}{\gboxend}

% 顶部导航栏（miniframes 主题：显示各节名称 + 帧圆点）
% 节名与圆点记录在 .nav 文件中，改动章节后需要编译两遍才会同步
\useoutertheme[subsection=false]{miniframes}
\setbeamercolor{section in head/foot}{fg=yellow, bg=maingreen}
\setbeamercolor{mini frame}{fg=yellow, bg=maingreen}
\setbeamercolor{mini frame in current subsection}{fg=yellow, bg=accentgreen}

% 顶部校徽栏：位于导航栏上方，宽度与帧标题栏一致（整页宽）
% 校训用的楷体：ctex 的 \kaishu 在本机指向未安装的 Adobe 楷体，这里直接指定系统楷体
\setCJKfamilyfont{kaiti}{KaiTi}
\setbeamercolor{logo bar}{fg=white, bg=darkgreen}
\addtobeamertemplate{headline}{%
  \begin{beamercolorbox}[wd=\paperwidth,ht=0.84cm,dp=0.20cm,leftskip=0.25cm,rightskip=0.25cm]{logo bar}%
    % -0.035cm：色条内容默认坐在基线上，会比栏的中线略高，下移这一点使其上下居中
    \raisebox{-0.035cm}{\includegraphics[height=0.72cm]{logo_white.png}}%
    % 右侧校训：字体为上面定义的楷体，也可改为 \insertsectionhead（当前节名）、\inserttitle（章节名）等
    % 0.375cm = 校徽中线（0.36cm）加汉字基线偏移，并同样下移 0.035cm 与校徽对齐
    \hfill\raisebox{\dimexpr 0.375cm - 0.5\height - 0.5\depth\relax}%
      {\CJKfamily{kaiti}\fontsize{12pt}{24pt}\selectfont 自强不息，止于至善}%
  \end{beamercolorbox}%
}{}

% % 设置边框 - 使用TikZ绘制圆角
% \setbeamertemplate{frametitle}{
%   \vspace*{0.2cm}
%   \begin{tikzpicture}
%     % 背景矩形宽度限制为标题内容宽度
%     \fill[lightgreen, rounded corners=5pt] (0,0) rectangle (\textwidth+0.6cm, 0.8cm);
%     \node[anchor=west, text width=\textwidth, xshift=0.3cm, yshift=0.4cm] {\usebeamerfont{frametitle}\insertframetitle};
%   \end{tikzpicture}
% }

% 设置项目符号
\setbeamertemplate{itemize item}{\small\raise0.5pt\hbox{\textcolor{accentgreen}{$\bullet$}}}
\setbeamertemplate{itemize subitem}{\small\raise0.3pt\hbox{\textcolor{accentgreen}{$\circ$}}}

% 设置页脚 - 使用TikZ绘制一条整页宽的色条，左侧页码、右侧姓名
\setbeamertemplate{footline}{
  \begin{tikzpicture}
    \fill[maingreen] (-1,0) rectangle (\paperwidth, 0.4cm);
    \node[yellow,anchor=east, xshift=0cm, yshift=0.2cm] {\insertframenumber/\inserttotalframenumber};
    % 原点在纸张左边缘右侧 1cm 处（色条从 x=-1cm 画起），故 \paperwidth-2cm
    % 对应纸张右边缘内缩 1cm 的位置，与左侧页码左右对称
    \node[yellow,anchor=east, xshift=\paperwidth-2cm, yshift=0.2cm] {\insertauthor};
  \end{tikzpicture}
}

% 关闭 beamer 默认的导航符号（页码与姓名已在页脚中自行绘制）
\setbeamertemplate{navigation symbols}{}

% 设置标题页面
\setbeamertemplate{title page}{
  \vspace{0.5cm}
  \begin{center}
    \begin{minipage}[c][0.68\paperheight][c]{\textwidth}
      {\usebeamerfont{title}\usebeamercolor[fg]{title}\large\inserttitle\par}
      \vspace{1.5cm}
      {\usebeamerfont{subtitle}\usebeamercolor[fg]{subtitle}\Large\insertsubtitle\par}
      \vspace{1.5cm}
      {\usebeamerfont{author}\usebeamercolor[fg]{author}\insertauthor\par}
      \vspace{0.5cm}
      {\usebeamerfont{institute}\usebeamercolor[fg]{institute}\insertinstitute\par}
      \vspace{0.5cm}
      {\usebeamerfont{date}\usebeamercolor[fg]{date}\insertdate\par}
    \end{minipage}
  \end{center}
}

% 设置块环境
\setbeamertemplate{blocks}[rounded][shadow=false]

% 设置目录
\setbeamertemplate{section in toc}{\inserttocsectionnumber.~\inserttocsection}


% ----------------------------------------------------------------------------
%  图片占位符（临时用；有图以后把那行换成 \includegraphics 即可）
%      \imgph{说明文字}          默认高 3.6cm
%      \imgph[6cm]{说明文字}     指定高度
%  会画一个绿色虚线框，框里写说明，方便对着 PDF 找位置。
\newcommand{\imgph}[2][3.6cm]{%
  \begin{center}%
  \begin{tikzpicture}%
    \node[draw=accentgreen, densely dashed, line width=1pt, rounded corners=4pt,
          fill=lightgreen, inner sep=8pt, align=center, text width=0.72\textwidth,
          minimum height=#1, font=\small]
      {\textcolor{accentgreen}{\bfseries 【图片占位】}\\[4pt] #2};%
  \end{tikzpicture}%
  \end{center}%
}

% ----------------------------------------------------------------------------
%  插图（图已就位后用这个；图放在 00-共享\图片\ 下，由 \graphicspath 找到）
%      \imgfig{文件名}                 默认宽 0.72\textwidth
%      \imgfig[0.55\textwidth]{文件名}  指定宽度
%      \imgfig[8cm]{文件名}            也可以直接给尺寸
%  图源与编译方法见 00-共享\图片\README.md。每张图是一页 PDF，只含图、无页边。
\newcommand{\imgfig}[2][0.72\textwidth]{%
  \begin{center}%
    \includegraphics[width=#1]{#2}%
  \end{center}%
}

% ----------------------------------------------------------------------------
%  可用方框环境速查（五类配色，均不编号，均支持 \begin{xxx}[补充名]）
%    定理类（深绿）  theorem 定理 / corollary 推论 / character 性质 /
%                    proposition 命题 / lemma 引理
%    定义类（海绿）  definition 定义 / axiom 公理
%    例题类（森林绿）example 例 / exercise 练习 / homework 习题
%    过程类（橄榄绿）proof 证明 / jie 解 / way1 法1 / way2 法2 /
%                    way3 法3 / way4 法4
%    注记类（青绿）  experience 经验总结 / zhu 注 / att 值得注意的是 /
%                    analyse 分析 / hint 提示
%  另有 beamer 自带的普通方框：\begin{block}{标题} ... \end{block}
% ----------------------------------------------------------------------------


% ==================== 每份课件只改这一段 ====================
\title{从数数开始的抽象代数生活}
\subtitle{第 13 课　它有什么用}
\author{姜博瀚}
\institute{厦门大学附属实验中学}
\date{\today}
% ============================================================

\begin{document}

\begin{frame}
  \titlepage
\end{frame}

% ============================================================================
\section{先回答第 1 课的问题}

\begin{frame}{费马小定理}
  \begin{theorem}[费马小定理]
    设 \(p\) 是素数，\(a\) 是不被 \(p\) 整除的整数，那么
    \[
      a^{\,p-1}\equiv1 \pmod p .
    \]
  \end{theorem}
  \pause
  \begin{center}
    {\large 工具已经磨好了。\pause 今天先拿它去解决一件旧事。}
  \end{center}
\end{frame}

\begin{frame}{黑板右边那行字}
  \begin{center}
    {\large 第 1 课上，我们在黑板最右边写下：}
    \[
      999999\div7=?\qquad\text{「第 13 课回来回答」}
    \]
  \end{center}
  \pause
  \begin{center}
    {\large 今天，把它划掉。}
  \end{center}
\end{frame}

\begin{frame}{\(999999\) 为什么能被 \(7\) 整除}
  第一步，改写：
  \[
    999999=1000000-1=10^{6}-1 .
  \]
  \pause
  第二步，用定理：\(7\) 是素数，\(10\) 不是 \(7\) 的倍数，取 \(p=7\)、\(a=10\)，于是
  \[
    10^{\,7-1}=10^{6}\equiv1 \pmod 7 .
  \]
  \pause
  第三步，代回去：
  \[
    10^{6}-1\equiv1-1=0 \pmod 7 .
  \]
  \pause
  \begin{center}
    {\large 余数是 \(0\)，所以 \(7\) \textbf{整除} \(999999\)。}\\[0.4em]
    {\large 证完。}
  \end{center}
\end{frame}

\begin{frame}{顺手把黑板右边那行字划掉}
  \begin{center}
    {\Large 第 1 课留下的问题：\quad \(\checkmark\)\quad 已解决}
  \end{center}
  \pause
  \begin{center}
    {\large 我们用了 \(13\) 节课，一次也没有用到计算器。}
  \end{center}
\end{frame}

% ============================================================================
\section{快速求大幂的余数}

\begin{frame}{\(2^{100}\) 除以 \(7\) 余几}
  不用算 \(2^{100}\)。看指数：
  \[
    100=6\times16+4 .
  \]
  \pause
  由费马小定理 \(2^{6}\equiv1\pmod 7\)，于是
  \[
    2^{100}=2^{6\times16+4}=(2^{6})^{16}\times2^{4}\equiv1^{16}\times16\equiv16\equiv2 \pmod 7 .
  \]
  \pause
  \begin{center}
    {\large \(2^{100}\) 除以 \(7\) 余 \(2\)。}
  \end{center}
\end{frame}

\begin{frame}{再练一道：\(3^{200}\) 除以 \(11\)}
  \[
    3^{10}\equiv1 \pmod{11}\quad(\text{费马小定理，}11\text{ 是素数}) .
  \]
  \pause
  \[
    200=10\times20 \quad\Longrightarrow\quad
    3^{200}=(3^{10})^{20}\equiv1^{20}=1 \pmod{11} .
  \]
  \pause
  \begin{center}
    {\large \(3^{200}\) 除以 \(11\) 余 \(1\)。}
  \end{center}
\end{frame}

\begin{frame}{用手指回答}
  \begin{itemize}
    \item \(2^{100}\) 除以 \(7\) 的余数是多少？
    \item 求 \(5^{2026}\) 除以 \(7\) 的余数，第一步应该看什么？
  \end{itemize}
  \pause
  \begin{block}{答案}
    \(2\)；\quad 先把 \(2026\) 除以 \(6\)（因为 \(5^6\equiv1\pmod 7\)），
    \(2026=6\times337+4\)，所以答案与 \(5^4\equiv4\) 相同。
  \end{block}
\end{frame}

% ============================================================================
\section{身份证上的那一位}

\begin{frame}{身份证最后一位，是怎么算出来的}
  身份证前 \(17\) 位，每一位乘一个固定的「权重」，加起来，再除以 \(11\) 取余。
  \begin{center}
  \begin{tabular}{@{}lc*{17}{c}@{}}
    \toprule
    位数 & \(1\) & \(2\) & \(3\) & \(4\) & \(5\) & \(6\) & \(7\) & \(8\) & \(9\) & \(10\) & \(11\) & \(12\) & \(13\) & \(14\) & \(15\) & \(16\) & \(17\) \\
    \midrule
    权重 & \(7\) & \(9\) & \(10\) & \(5\) & \(8\) & \(4\) & \(2\) & \(1\) & \(6\) & \(3\) & \(7\) & \(9\) & \(10\) & \(5\) & \(8\) & \(4\) & \(2\) \\
    \bottomrule
  \end{tabular}
  \end{center}
  \pause
  算出来的余数，按一张表换成最后一位：
  \begin{center}
  \begin{tabular}{@{}l*{11}{c}@{}}
    \toprule
    余数 & \(0\) & \(1\) & \(2\) & \(3\) & \(4\) & \(5\) & \(6\) & \(7\) & \(8\) & \(9\) & \(10\) \\
    \midrule
    校验码 & \(1\) & \(0\) & X & \(9\) & \(8\) & \(7\) & \(6\) & \(5\) & \(4\) & \(3\) & \(2\) \\
    \bottomrule
  \end{tabular}
  \end{center}
\end{frame}

\begin{frame}{算一个给你看}
  取一个例子编号的前 \(17\) 位：
  \[
    3\ 5\ 0\ 1\ 0\ 2\ 1\ 9\ 9\ 0\ 0\ 3\ 0\ 7\ 0\ 0\ 1
  \]
  每一位乘上对应的权重，加起来：
  \[
    21+45+0+5+0+8+2+9+54+0+0+27+0+35+0+0+2=208 .
  \]
  \pause
  \[
    208=11\times18+10 \quad\Longrightarrow\quad \text{余数 }10 \;\longrightarrow\; \text{校验码 }2 .
  \]
  \pause
  \begin{center}
    所以完整的编号是 \quad \(\mathbf{350102199003070012}\)。
  \end{center}
\end{frame}

\begin{frame}{为什么这一位能防错}
  \begin{itemize}
    \item 前 \(17\) 位里改掉任意一位，加权和几乎一定会变；
    \item 余数变了，校验码就对不上，机器立刻知道「这个号不对」。
  \end{itemize}
  \pause
  \begin{att}
    这背后全是「取余」——和我们课上算的模 \(11\) 是同一件事。
  \end{att}
\end{frame}

% ============================================================================
\section{RSA：一个极小的例子}

\begin{frame}{怎么「锁」，怎么「开」}
  取两个素数 \(p=3\)、\(q=11\)，令 \(n=3\times11=33\)。
  \pause
  \begin{block}{公钥（谁都能拿）}
    \(n=33\) 和加密钥匙 \(e=3\)。\quad 加密规则：把数 \(x\) 变成 \(x^{3}\) 除以 \(33\) 的余数。
  \end{block}
  \pause
  \begin{block}{私钥（只有收信人有）}
    解密钥匙 \(d=7\)。\quad 解密规则：把收到的数 \(y\) 变成 \(y^{7}\) 除以 \(33\) 的余数。
  \end{block}
\end{frame}

\begin{frame}{把 5 加密，再解开}
  \begin{center}
  \begin{tabular}{@{}lcc@{}}
    \toprule
    步骤 & 算法 & 结果 \\
    \midrule
    要送出去的数 & \(x\) & \(5\) \\
    加密 & \(5^{3}=125\)，再除以 \(33\) 取余 & \(26\) \\
    传过去 & & \(26\) \\
    解密 & \(26^{7}\) 除以 \(33\) 取余 & \(5\) \\
    \bottomrule
  \end{tabular}
  \end{center}
  \pause
  \begin{center}
    {\large 收到的又是 \(5\)。锁上、打开，都对。}
  \end{center}
\end{frame}

\begin{frame}{这门课为什么讲得到它}
  \begin{itemize}
    \item 「锁」和「开」用的都是\textbf{取余}；
    \item 为什么 \(d=7\) 就对、别人猜不出来？——这里藏着费马小定理的推广；
    \item 真正在用的 \(n\) 有几百度长，靠的是「拆开容易、合起来难」。
  \end{itemize}
  \pause
  \begin{center}
    {\large 手机里、网页上，每一次「https」用的都是这类\textbf{绕圈的算术}。}
  \end{center}
\end{frame}

\begin{frame}{这节课要带走的三句话}
  \begin{itemize}
    \item 费马小定理能回答 \(999999\div7\)：\(10^{6}\equiv1\pmod 7\)，所以整除。
    \item 它还能把巨大的幂快速地「缩」小：\(2^{100}\equiv2\pmod 7\)。
    \item 身份证校验码、手机里的加密，用的都是这套「取余」的算术。
  \end{itemize}
  \pause
  \begin{center}
    {\large\textcolor{textgreen}{\bfseries 下节课：把整门课串成一条线。}}
  \end{center}
\end{frame}

% ---------------- 结尾页 ----------------
\begin{frame}
  \begin{center}
    \vspace{1.6cm}
    {\Large\textcolor{textgreen}{\bfseries 下课}}

    \vspace{1.0cm}
    {\large\textcolor{accentgreen}{下节课：复习、总结与展望}}
  \end{center}
\end{frame}

\end{document}